\addtocontents{toc}{\vspace{-1.0em}}
\mysection{Обобщени мрежи: Теоретични основи и софтуер за симулации}


\subsection{Бележки върху теорията на обобщените мрежи}
\subsubsection{За понятието обобщена мрежа}

\hspace*{2.5em}Обобщените мрежи (ОМ) са представени за пръв път през 1982 г. в доклад на международна конференция, като текстът е публикуван през 1984 г. \cite{XY20}. Те се появяват като естествено продължение на усилията за моделиране на паралелни процеси, започнали с изследванията на \textit{C.~A.~Petri} от 1962 г., довели до формулирането на мрежите на Петри (МП) \cite{XY1}. В класическия си вид МП са двуделен ориентиран граф с два типа върхове: \emph{преходи} (дискретни събития) и \emph{позиции} (условия), между които се движат \emph{ядра} (\textit{tokens}). Позициите могат да бъдат входни, изходни или вътрешни; в зависимост от ориентацията на ребрата те действат като предусловия или постусловия на преходите, а преместването на ядрата се разрешава само при реализируем (разрешен) преход.

Насочването към паралелизма води до множество разширения на МП (над двайсет до 1982 г.), които добавят: (i) времеви аспекти: моменти на преместване на ядрата и продължителност на преходите; (ii) оцветяване на ядрата за различимост; (iii) охраняващи условия (предикати) върху групи ядра за синхронизирани преминавания; и др. Значима част от тях са \emph{съществени разширения}, т.е. описват процеси, недостижими за базовата МП. В този контекст се вписват и конкретни семейства като E-мрежите на \textit{G.~Nutt} \cite{XY2} и времевите варианти на МП (\textit{Ramchandani} \cite{XY3}, \textit{Merlin} \cite{XY4}, \textit{Crespi-Reghizzi} и \textit{Mandrioli} \cite{XY5}).

На този фон ОМ предлагат по-богата рамка за спецификация на едновременност, синхронизация и условни зависимости, запазвайки интуитивната графова семантика на МП, но разширявайки изразителната мощ за моделиране на сложни паралелни процеси. Подробни разглеждания на ОМ са дадени в \cite{Z3,Z1,BAM, AtanassovSotirova2017}, а малка част от техните приложения са разгледани в \cite{Orozova2011,Georgiev1998,StefanovaPavlova2001,Aladjov2002,SgurevAIHistory2021, ZotevaKrawczakDMSurvey2017}.

\subsubsection{Неформално описание на означенията в обобщените мрежи}

\hspace*{2.5em}Тук обобщаваме основните означения и правила в ОМ, като се позоваваме на примерната схема на \textbf{Фигура~\ref{fig:generalized-net-example}} и общия вид на прехода на \textbf{Фигура~\ref{fig:gn-transition}}. 
Нейните \it позиции \rm са показани с \hspace{-2mm}
\unitlength 1.0mm
\linethickness{0.4pt}
\begin{picture}(10.00,4.00)
\put(3.0,2.0){\circle{4.0}}
\end{picture}
\hspace{-6mm}, а всеки елемент от вида на \textbf{Фигура~\ref{fig:gn-transition}} е \em преход\rm. Преходът съдържа условия, които са означени с 
\unitlength 1.00mm
\begin{picture}(2.00,6.00)
\put(1.,-0.50){\line(0,1){5.00}}
\put(1.03,4.00){\vector(0,-1){.2}}
\end{picture}\,. 
Позиции и преходи се свързват с \it ребра\rm\ (стрелки), които определят възможните направления на движение.

В някои позиции се намират \it ядра\rm\ (точки), аналогично на класическите мрежи на Петри~\cite{Q148}. Входящото ядро носи \textit{начални характеристики}, а при преминаване през мрежата може да получи допълнителни характеристики. Тъй като първоначалните характеристики могат да останат валидни, за ядрото се натрупва \textit{история}. От гледна точка на реализацията, това множество от характеристики може да се възприема като асоциативен списък „име~$\to$~стойност“. Позицията може да има \it характеристична функция\rm, която задава нови стойности на характеристиките на ядрата, влизащи в нея.

Свързването в ОМ е опростено по отношение на степените: всяка позиция е свързана с поне едно ребро и допуска най-много едно входно и едно изходно ребро. Без входно ребро позицията е \it входна\rm, а без изходно -- \it изходна\rm.

В примера на \textbf{Фигура~\ref{fig:generalized-net-example}} $l_{1}$, $l_{2}$ и $l_{5}$ са входни, а $l_{6}$ и $l_{8}$ -- изходни. Всеки преход има поне по една входна и изходна позиция ($m,n\!\ge\!1$); възможно е при даден преход една и съща позиция да играе и двете роли \cite{Z1}. Когато ядро достигне входна позиция на преход, то става \em потенциално активно\rm; ако са налице условията за прехвърляне към изходните позиции, преходът е \it активен\rm.

Условията към преходите не се описват с единичен предикат, а чрез \it индексирана матрица\rm\ (ИМ) -- матрица от предикати, която едновременно специфицира кога е възможно движение по различните направления и определя капацитетите на съответните ребра. Така ИМ осигурява по-фин контрол спрямо стандартните мрежи на Петри \cite{Z1}.

Времето в дефиницията на ОМ е \it дискретно\rm: разглеждат се последователни стъпки, номерирани с целочислена променлива, която нараства с единица. Допустимо е паралелно да се фиксира \it абсолютна времева скала\rm\ с начален момент и хоризонт на функциониране (например начало на 1.10.1964~г. и продължителност 50 години или до края на XXI век) \cite{Z1}. При множество свързани процеси може да се използват отделни скали за всеки или обща за синхронизация.

ОМ са достатъчно изразителни, за да опишат както разширенията на мрежите на Петри, така и модели с изчислителна мощ от класа на машината на Тюринг (в частност и крайни автомати) \cite{Z1}. Известните към момента разширения на ОМ са \it консервативни\rm: за всяко от тях съществува стандартна ОМ, която еквивалентно описва функционирането и резултата му \cite{Z1,BAM}.

\vspace{0.2cm}

\begin{figure}[htbp]
	\centering
	\unitlength 1mm
	\linethickness{0.4pt}
	
	\begin{picture}(75,76)(4,-4)
		% Позиции (Places) 
		\put(7,55){\makebox(0,0)[cb]{$l_{1}$}}
		\put(7,50){\circle{6.0}}
		\put(7,33){\makebox(0,0)[cb]{$l_{2}$}}
		\put(7,28){\circle{6.0}}
		
		\put(25,11){\makebox(0,0)[cb]{$l_{5}$}}
		\put(25,6){\circle{6.0}}
		
		\put(25,55){\makebox(0,0)[cb]{$l_{3}$}}
		\put(25,50){\circle{6.0}}
		\put(25,33){\makebox(0,0)[cb]{$l_{4}$}}
		\put(25,28){\circle{6.0}}
		
		\put(43,55){\makebox(0,0)[cb]{$l_{6}$}}
		\put(43,50){\circle{6.0}}
		\put(43,33){\makebox(0,0)[cb]{$l_{7}$}}
		\put(43,28){\circle{6.0}}
		
		\put(61,55){\makebox(0,0)[cb]{$l_{8}$}}
		\put(61,50){\circle{6.0}}
		\put(61,33){\makebox(0,0)[cb]{$l_{9}$}}
		\put(61,28){\circle{6.0}}
		\put(61,11){\makebox(0,0)[cb]{$l_{10}$}}
		\put(61,6){\circle{6.0}}
		
		% Преходи (Transitions)
		\put(16,68){\makebox(0,0)[cb]{$Z_{1}$}}
		\put(16,65){\vector(0,-1){0.27}}
		\put(16,2){\line(0,1){63}}
		
		\put(34,68){\makebox(0,0)[cb]{$Z_{2}$}}
		\put(34,65){\vector(0,-1){0.27}}
		\put(34,2){\line(0,1){63}}
		
		\put(52,68){\makebox(0,0)[cb]{$Z_{3}$}}
		\put(52,65){\vector(0,-1){0.27}}
		\put(52,2){\line(0,1){63}}
		
		% Дъги (Arcs)
		\put(10,50){\vector(1,0){6.0}}
		\put(16,50){\vector(1,0){6.0}}
		\put(10,28){\vector(1,0){6.0}}
		\put(16,28){\vector(1,0){6.0}}
		
		\put(28,50){\vector(1,0){6.0}}
		\put(34,50){\vector(1,0){6.0}}
		\put(28,28){\vector(1,0){6.0}}
		\put(34,28){\vector(1,0){6.0}}
		\put(28,6){\vector(1,0){6.0}}
		
		\put(52,50){\vector(1,0){6.0}}
		\put(46,28){\vector(1,0){6.0}}
		\put(52,28){\vector(1,0){6.0}}
		\put(52,6){\vector(1,0){6.0}}
		
		% Обратни връзки (Feedback)
		% l10 -> z3
		\put(64,6){\line(1,0){6.0}}
		\put(70,6){\line(0,-1){6}}
		\put(70,0){\line(-1,0){24}}
		\put(46,0){\line(0,1){6}}
		\put(46,6){\vector(1,0){6.0}}
		
		% l9 -> z1
		\put(64,28){\line(1,0){12.0}}
		\put(76,28){\line(0,-1){30}}
		\put(76,-2){\line(-1,0){66}}
		\put(10,-2){\line(0,1){8}}
		\put(10,6){\vector(1,0){6.0}}
		
		% Вътрешна зависимост
		\put(28,28){\line(1,0){6.0}}
		\put(34,28){\line(0,1){22.0}}
		\put(34,50){\line(1,0){3.0}}
		\put(37,50){\vector(1,0){3.0}}
	\end{picture}
	\caption{Пример за обобщена мрежа}
	\label{fig:generalized-net-example}
\end{figure}

Във \textbf{Фигура~\ref{fig:generalized-net-example}} $l_1 \dots l_{10}$ са позиции, а $Z_1$, $Z_2$ и $Z_3$ са преходи. А във \textbf{Фигура~\ref{fig:gn-transition}} е даден общ вид на преход в ОМ.

\begin{figure}[htbp]
	\centering
	\unitlength 1mm
	\linethickness{0.4pt}
	
	\begin{picture}(64,60.00)(20,0)
		% Лява част (входни позиции)
		\put(24.16,40.16){\makebox(0,0)[cb]{$\vdots$}}
		\put(24.16,20.16){\makebox(0,0)[cb]{$\vdots$}}
		\put(50.16,40.16){\makebox(0,0)[cb]{$\vdots$}}
		\put(50.16,20.16){\makebox(0,0)[cb]{$\vdots$}}
		
		\put(24.16,50.16){\makebox(0,0)[cb]{$l'_{1}$}}
		\put(32.16,50.16){\circle{6.32}}
		\put(35.32,50.00){\vector(1,0){16.68}}
		
		\put(24.16,30.16){\makebox(0,0)[cb]{$l'_{i}$}}
		\put(32.16,30.16){\circle{6.32}}
		\put(35.32,30.00){\vector(1,0){16.68}}
		
		\put(24.16,10.16){\makebox(0,0)[cb]{$l'_{m}$}}
		\put(32.16,10.16){\circle{6.32}}
		\put(35.32,10.00){\vector(1,0){16.68}}
		
		% Централна вертикална линия Z
		\put(52.00,57.00){\makebox(0,0)[cb]{$Z$}}
		\put(52.00,53.00){\vector(0,-1){0.1}}
		\put(52.00,2.00){\line(0,1){51.00}}
		
		% Дясна част (изходни позиции)
		\put(80.16,36.16){\makebox(0,0)[cb]{$\vdots$}}
		\put(80.16,16.16){\makebox(0,0)[cb]{$\vdots$}}
		\put(54.16,36.16){\makebox(0,0)[cb]{$\vdots$}}
		\put(54.16,16.16){\makebox(0,0)[cb]{$\vdots$}}
		
		\put(80.00,46.16){\makebox(0,0)[cb]{$l''_{1}$}}
		\put(72.00,46.16){\circle{6.32}}
		\put(52.00,46.00){\vector(1,0){16.84}}
		
		\put(80.00,26.16){\makebox(0,0)[cb]{$l''_{j}$}}
		\put(72.00,26.16){\circle{6.32}}
		\put(52.00,26.00){\vector(1,0){16.84}}
		
		\put(80.00,6.16){\makebox(0,0)[cb]{$l''_{n}$}}
		\put(72.00,6.16){\circle{6.32}}
		\put(52.00,6.00){\vector(1,0){16.84}}
	\end{picture}
	\caption{Общ вид на преход в ОМ}
	\label{fig:gn-transition}
\end{figure}


\subsubsection{Формална дефиниция на обобщените мрежи}

\hspace*{2.5em}Въз основа на изложеното неформално описание, настоящият раздел представя строго математическо дефиниране на компонентите на ОМ, следвайки \cite{Z3}. Формалният апарат се базира на теорията на множествата и апарата на ИМ, като дефинира структурата, времевите параметри и динамиката на ядрата. За нуждите на изложението се въвеждат следните базови означения:

\begin{itemize}
	
	\item ${\cal N} = \lbrace 0, 1, 2, \dots \rbrace \cup \lbrace \infty \rbrace$;
	
	\item Нека $X=\langle x_1,x_2,\dots,x_n\rangle$ е $n$-мерен вектор, където
	$n \in {\cal N}$ и $n \ge 1$.
	С $\mbox{pr}_{i}X$ означаваме $i$-тата координатна проекция на $X$, т.е. тя извежда
	$i$-тата координата на вектора:

	\begin{equation}
	\mbox{pr}_{i}(X) = x_i,
	\qquad 1 \le i \le n .
	\label{eq:projection-single}
	\end{equation}

	По-общо, за множество индекси $I=\{i_1,i_2,\dots,i_k\}\subseteq \{1,2,\dots,n\}$,
	избрани така, че
	$1 \le i_1 < i_2 < \dots < i_k \le n$ и $1 \le k \le n$,
	дефинираме координатната проекция върху избраните координати
	$\mbox{pr}_{i_1,i_2,\dots,i_k}(X)$ чрез:

	\begin{equation}
	\mbox{pr}_{i_1,i_2,\dots,i_k}(X)
	=
	\langle \mbox{pr}_{i_1}(X), \mbox{pr}_{i_2}(X), \dots, \mbox{pr}_{i_k}(X) \rangle
	=
	\langle x_{i_1}, x_{i_2}, \dots, x_{i_k} \rangle .
	\label{eq:projection-subvector}
	\end{equation}

	Както е показано във формула \eqref{eq:projection-subvector}, резултатът е $k$-мерен вектор,
	получен чрез подреждане на избраните координати на $X$.
	Например, за $X=\langle 5,4,3,2,1\rangle$ получаваме
	$\mbox{pr}_{2,3,5}(X)=\langle 4,3,1\rangle$.
	
\end{itemize}

Формално всеки преход се описва чрез наредената седморка: \index{преход}
$$ Z = \langle L', L'', t_{1}, t_{2}, r, M, \kv \rangle, $$ 

%here i put newpage command
\newpage
където:

{\bf (а)} $L'$ и $L''$ са крайни непразни множества от съответно входни и изходни позиции на прехода. За прехода, показан на \textbf{Фигура~\ref{fig:gn-transition}}, те са:
$$ L' = \{ l'_{1}, l'_{2}, \dots, l'_{m} \} \quad \text{и} \quad L'' = \{ l''_{1}, l''_{2}, \dots, l''_{n} \}; $$

{\bf (б)} $t_{1}$ е моментът във времето, когато преходът може да се активира, т.е. времето, когато всички условия, необходими за активацията, са изпълнени, и преходът става възможен за извършване;

{\bf (в)} $t_{2}$ е текущата продължителност на активното състояние на прехода, имаща отношение към динамиката и промените в състоянието на системата;

{\bf (г)} $r$ е ИМ, която задава условията на прехода и има вида:
$$r = \begin{array}{c|c}
& l''_{1} \mbox{     } \dots \mbox{     }   l''_{j}\mbox{     }   \dots \mbox{     }  l''_{n}  \\ \hline
  l'_{1} & \\
  \vdots &   r_{i,j} \\
  l'_{i} &  (r_{i,j} \spa \spa - \mbox{предикат}, \\
  \vdots &   1 \leq i \leq m, 1 \leq j \leq n) \\
  l'_{m} & \\
\end{array} \spa \spa $$
където $r_{i,j}$ е предикатът, който, ако в текущия момент има стойност „истина“, разрешава преминаването на ядро от $i$-тата входна позиция в $j$-тата изходна позиция. Ако обаче предикатът не е изпълнен (т.е. има стойност „неистина“), преходът от тази позиция в друга става невъзможен.

{\bf (д)} $M$ е ИМ, която задава капацитетите на ребрата на прехода, и има вида:
$$M = \begin{array}{c|c}
& l''_{1} \mbox{     } \dots \mbox{     }  l''_{j} \mbox{     } \dots\mbox{     }  l''_{n}  \\ \hline
  l'_{1} & \\
  \vdots &   m_{i,j} \\
  l'_{i} &  (m_{i,j} \geq 0  - \mbox{естествено число или $\infty$}, \\
  \vdots &   1 \leq i \leq m, 1 \leq j \leq n) \\
  l'_{m} & \\
\end{array} \spa \spa $$

Тези капацитети $m_{i,j}$ определят колко „ядра“ могат да преминават по реброто от входната позиция $l'_i$ към изходната позиция $l''_j$. Например, ако $m_{i,j} = 3$, това означава, че през това ребро могат да преминат до 3 ядра, а ако $m_{i,j} = \infty$, то няма ограничение за броя на ядрата.

{\bf (е)} $\kv$ е типът на прехода, който се представя като булев израз. Аргументите на израза са идентификаторите на входните позиции на прехода, свързани с логическите операции $\wedge$ (и) и $\vee$ (или). За всяко подмножество от позиции $\{ l_{i_{1}}, \dots, l_{i_{u}} \} \subset L'$ са дефинирани следните условия:

$$ \begin{array}{ll}
     \wedge(l_{i_{1}}, l_{i_{2}}, \dots, l_{i_{u}}) & \mbox{изисква всяка от позициите } \\
     & \mbox{да съдържа поне едно ядро;} \\
    \vee(l_{i_{1}}, l_{i_{2}}, \dots, l_{i_{u}}) & \mbox{изисква наличието на поне} \\
    & \mbox{едно ядро в някоя от} \\
     & \mbox{изброените позиции.}
\end{array} $$


% Когато настъпи моментът за активиране на прехода и ако булевият израз е истина, преходът се активира. В противен случай преходът не се активира и състоянието на системата не се променя.

Преходът се активира в момента $t_1$, при условие че вярностната стойност на булевия израз е „истина“. В противен случай преходът остава пасивен и състоянието на системата не се променя.

Наредената четворка
$$ E = \langle \langle A, \pi_{A}, \pi_{L}, c, f, \theta_{1}, \theta_{2} \rangle, 
\langle K, \pi_{K}, \theta_{K} \rangle, \langle T, t^{0}, t^{*} \rangle, 
\langle X, \Phi, b \rangle \rangle $$ 
се нарича \em обобщена мрежа \rm (ОМ), ако: \index{обобщена мрежа (ОМ)}

{\bf (а)} $A$ е множеството на преходите на ОМ. Това множество съдържа всички възможни преходи, които могат да се случат в рамките на мрежата, като всеки преход представлява операция върху състоянията на позициите.

{\bf (б)} $\pi_{A}$ е функция, която задава приоритетите на преходите. Тоест, $\pi_{A}: A \rightarrow {\cal N}$, като за всеки преход в множеството $A$ се определя неговият приоритет, който е елемент от множеството на естествените числа ${\cal N}$.

{\bf (в)} $\pi_{L}$ е функция, която задава приоритетите на позициите в ОМ. Тоест, $\pi_{L}: L \rightarrow {\cal N}$, където:
$$L = \mbox{\mbox{pr}}_{1} A \cup \mbox{\mbox{pr}}_{2} A$$
е множеството на всички позиции на ОМ. Позициите в ОМ се разпределят по приоритети, като всяка позиция получава приоритет, определен от функцията $\pi_L$.

{\bf (г)} $c$ е функция, която задава капацитетите на позициите, т.е. $c: L \rightarrow {\cal N}$. Тази функция показва колко „ядра“ или елементи могат да бъдат свързани с всяка позиция в мрежата. Например, ако $c(l) = 5$, това означава, че на позицията $l$ могат да бъдат свързани до 5 ядра.

{\bf (д)} $f$ е функция, определяща вярностната стойност на предикатите в условията на преходите. В стандартните ОМ тя приема стойности от множествата $\{ \text{false, true} \}$ или $\{ 0, 1 \}$, които индикират дали дадено логическо условие за активиране на прехода е изпълнено.

{\bf (е)} $\theta_{1}$ е функция, определяща следващия момент на активиране на прехода: $\theta_{1}(t) = t'$, където $t < t'$ и $t, t' \in [T, T + t^{*}]$. Стойността ѝ се изчислява при стартиране на мрежата и след всяко преустановяване на активността на прехода.

{\bf (ж)} $\theta_{2}$ е функция, определяща продължителността на активното състояние на прехода $Z$: $\theta_{2}(t) = t'$, където $t \in [T, T + t^{*}]$ и $t' > 0$. Стойността ѝ се изчислява в момента на активиране на прехода и дефинира неговия времеви обхват.

{\bf (з)}  $K$ е множеството на всички ядра на ОМ, което включва всички елементи, които взаимодействат с мрежата и участват в нейното функциониране. Ядрата могат да бъдат различни физически или логически обекти, които се движат или взаимодействат със състоянията в ОМ.

{\bf (и)}  $\pi_{K}$ е функция, задаваща приоритетите на ядрата, т.е. $\pi_{K}: K \rightarrow {\cal N}$, където всяко ядро получава числов приоритет, който показва неговото значение или важност за мрежата. Тази функция може да се използва за управление на реда на обработка на ядрата в ОМ.

{\bf (й)}  $\theta_{K}$ е функция, задаваща момента от време, в който дадено ядро може да влезе в ОМ, т.е. $\theta_{K}(\alpha) = t$, където $\alpha \in K$ и $t \in [T, T + t^{*}]$. Тази функция контролира времето, когато ядрата стават активни и започват да участват в процесите в ОМ.

{\bf (к)} $T$ е моментът, определен по фиксирана времева скала, в който ОМ започва да функционира; той дефинира началото на времевия интервал на работа на мрежата и нейните процеси.

{\bf (л)}  $t^{0}$ е елементарната времева стъпка, с която нараства времето във фиксираната времева скала. Това е минималният интервал от време, който може да бъде измерен или обработен в ОМ.

{\bf (м)} $t^{*}$ е общата продължителност на функциониране на мрежата; заедно с началния момент $T$, тя дефинира времевия интервал $[T, T + t^{*}]$, в който ОМ изпълнява своите процеси.

{\bf (н)} $X$ е функция, дефинираща началната характеристика $x^{\alpha}_{0}$ на всяко ядро $\alpha$ при неговото постъпване в мрежата.

{\bf (о)} $\Phi$ е характеристична функция, определяща новата характеристика на всяко ядро $\alpha$ при преминаването му в следваща позиция на мрежата, като текущата му характеристика се означава с $x^{\alpha}_{cu}$.

{\bf (п)} $b$ е функция, определяща максималния брой съхранявани характеристики за всяко ядро, $b: K \rightarrow \mathbb{N} \cup \{ \infty \}$. Стойността $b(\alpha)$ дефинира капацитета на паметта за ядрото $\alpha$, като:
\begin{itemize}[leftmargin=3.9em, nosep]
	\item при $b(\alpha) = 1$ се съхранява само текущата характеристика;
	\item при $b(\alpha) = \infty$ се пази пълната история на характеристиките;
	\item при $b(\alpha) = k < \infty$ се съхраняват началната и последните $k$ характеристики.
\end{itemize}

\vspace{0.3cm}

Компонентите на ОМ могат да бъдат класифицирани според тяхната функционална роля в модела:

\vspace{-0.3cm}

\begin{itemize}
	\item \textbf{Статичната структура} се определя от елементите на $\text{pr}_{1,2,6,7} A$ (входни/изходни позиции, капацитети на ребрата и тип на преходите), както и от функциите за приоритет $\pi_{A}, \pi_{L}$ и капацитети $c$;
	\item \textbf{Динамичната природа} се обуславя от ядрата, условията на преходите и функциите $f$ и $\pi_{K}$;
	\item \textbf{Времевата същност} се изразява чрез параметрите $T, t^{0}, t^{*}$, елементите на $\text{pr}_{3,4} A$ и функциите $\theta_{1}, \theta_{2}, \theta_{K}$;
	\item \textbf{Паметта на ОМ} се формира от компонентите $\Phi, X$ и $b$, управляващи характеристиките на ядрата.
\end{itemize}

Дадена ОМ се нарича \it празна ОМ \rm, ако $A = \varnothing$ или $K = \varnothing$.

Нека $\Sigma$ е класът на всички ОМ. Той е същински клас в смисъла на теорията на множествата \cite{TM1}.

Удобно е да означим с $Q^I$ и $Q^O$ множествата на входните и на изходните позиции на ОМ.


\subsubsection[Алгоритми за движение на ядрата в преход и в обобщена мрежа]{Алгоритми за движение на ядрата в преход и в обобщена мрежа}\label{tokenMove}

\hspace*{2.5em}Тук се представя разширена версия на най-общия от познатите досега алгоритми за функциониране на преход в ОМ \cite{Z1, Z3, BAM, AtanassovSotirova2017}. В базовата версия на алгоритъма, дефинирана в \cite{BAM}, не е предвидена възможност за сливане на входящо ядро със съществуващо такова в изходните позиции. Поради това, при достигане на максималния капацитет на дадена позиция, тя блокира приемането на нови ядра в рамките на текущата времева стъпка.

В предлагания вариант се допуска ядрото с най-висок приоритет от входна позиция (при удовлетворен съответен предикат) да се слее с някое от наличните ядра в запълнена до капацитета си изходна позиция. По този начин броят на обектите в нея остава непроменен и капацитетът \'{и} не се надвишава.

Важно изискване към моделите, илюстрирано в \cite{Z1, Z3}, е предикатите да не зависят от бъдещи събития, които потенциално биха настъпили в мрежата.

\vspace{0.5em}

\textbf{Алгоритъм А за функциониране на преход в ОМ} 

Приема се, че началната характеристика на дадено ядро $\alpha$ съдържа информация за това с кои 
други ядра може да се слее. Удобно е като втора компонента да се използва символ 
$S \in \{Y_S, N_S\}$, където $Y_S$ означава, че ядрото подлежи на разцепване, а $N_S$ -- че не подлежи.  

Например, началната характеристика на ядрото $\alpha$ може да се формализира като:
\begin{equation}
	x^{\alpha}_0 = \langle \{\beta_1, \dots, \beta_k\}, Y_S, x^{\alpha,*}_0 \rangle 
	\label{eq:alpha-initial-characteristic}
\end{equation}

Както е показано във формула~\eqref{eq:alpha-initial-characteristic}, това означава, че $\alpha$ може да се разцепва, 
да се слива с елементите на множеството $\{\beta_1, \dots, \beta_k\}$, а $x^{\alpha,*}_0$ съдържа остатъчната информация, 
която се съхранява в началната характеристика на ядрото.

\vspace{0.5cm}

{\setlength{\parskip}{0pt}  % Local override for this paragraph
\textbf{Алгоритъмът включва следните 14 стъпки:}
}

\vspace{-0.4cm}

\begin{itemize}
  \item (А01) Входните и изходните позиции на прехода се подреждат спрямо приоритета си.

  \item (А02) Ядрата във входните позиции се сортират по приоритет. Всяка позиция се разделя на 
  две групи ядра: (1) група от ядра, които могат да преминат в изходните позиции в текущата стъпка; 
  (2) първоначално празна група, където ще попаднат останалите ядра.

  \item (А03) Инициализира се празна ИМ $R$, която е аналог на ИМ $r$, дефинирана с предикатите 
  на условието за преход. За всички елементи $R_{i,j}$ в $R$:
  \begin{itemize}[itemsep=0.5em]
    \item (А03а) Присвоява се $0$ на $R_{i,j}$, ако $i$-тата входна позиция е празна, тоест в нея 
    няма ядра, годни да се придвижат към изходна позиция.
    \item (А03б) Присвоява се $0$ на $R_{i,j}$, ако $j$-тата изходна позиция вече няма 
    свободно място (пълна е).
    \item (А03в) Присвоява се $0$ на $R_{i,j}$, когато $r_{i,j} = \textit{false}$ или $m_{i,j} = 0$, 
    т.е. когато предикатът не е изпълнен или капацитетът на реброто между $i$-тата входна 
    и $j$-тата изходна позиция е изчерпан.
    \item (А03г) На всички останали $R_{i,j}$, за които $r_{i,j} = \textit{true}$, се задава стойност $1$.
  \end{itemize}

\item (А04) Следва се поредицата от входни позиции според техния приоритет, тръгвайки от тази с най-висок приоритет, където има поне едно ядро, което досега не е преминало към изходна позиция в текущия времеви интервал. Изпълняват се следните подстъпки:
  \begin{itemize}[itemsep=0.5em]
    \item (А04а) Проверява се актуалната стойност на $R_{i,j}$.
      \begin{itemize}
        \item Ако $R_{i,j}$ \emph{не е дефинирано}, се преминава към (А04б).
        \item Ако $R_{i,j}=1$, се преминава към (А04в).
        \item Ако $R_{i,j}=0$, се проверява дали в съответната изходна позиция има ядро, с което текущото ядро би могло да се слее. При положителен отговор се преминава към (А04в), иначе към (А04д).
      \end{itemize}
    \item (А04б) Проверява се предикатът $r_{i,j}$ в $r$. Ако е \emph{истина}, се преминава към (А04в), в противен случай -- към (А04д).
    \item (А04в) Придвижва се текущото ядро към $j$-тата изходна позиция. Ако в нея има подходящо ядро, се извършва сливане. Актуализират се съответните заетости/капацитети (в т.ч. $m_{i,j}$ и състоянието на изходната позиция). След преместването се преминава към (А04г).
    \item (А04г) Ако ядрото не се слее с друго ядро или след сливане не съществуват повече ядра в другите входни позиции, с които то би могло да се обедини, то получава следващата си характеристика. В случай че има още ядра за потенциално сливане, то „изчаква“ (вече преместено в изходната позиция), докато се извърши последното възможно сливане, и едва тогава получава нова характеристика -- в текущия или в заключителния такт от функционирането на прехода.
    \item (А04д) Ако ядрото не може да се разцепва или вече са проверени всички предикати от съответния ред в $r$, се преминава към (А05). Иначе се връща към (А04а) за следващ допустим $(i,j)$-кандидат според приоритета.
  \end{itemize}


  \item (А05) Ако ядрото не може да се придвижи на тази стъпка, то се добавя към втората група 
  ядра в текущата входна позиция.

  \item (А06) За всички изходни позиции, в които е постъпило ядро, без да се слее с друго ядро, броят на ядрата се увеличава с $1$.

  \item (А07) Броят на ядрата в онези входни позиции, откъдето ядро е било преместено към изходна 
  позиция, намалява с 1. Ако входната позиция остане без ядра, всички елементи в съответния ред на $R$ приемат стойност нула.

  \item (А08) Капацитетите на всички ребра, по които е преминало ядро, се редуцират с $1$.

  \item (А09) Ако има все още входни позиции с по-нисък приоритет, от които ядро не е преминало 
  към изходните позиции, алгоритъмът се връща на (А04). В противен случай се преминава към (А10).

  \item (А10) Текущото време се увеличава с $t^0$.

  \item (А11) Ако текущото време е достигнало или надхвърлило $t_1 + t_2$, се преминава към стъпка (А12), в противен случай се преминава към стъпка (А04).

  \item (А12) Работата на прехода приключва.
\end{itemize}


\textbf{Алгоритъм Б за функциониране на ОМ}

За описанието на алгоритъма се въвежда понятието \emph{абстрактен преход}, което представлява обединението на всички преходи, активни в даден времеви момент.

\begin{itemize}
  \item (Б01) Всички ядра $\alpha$, за които $\theta_{K}(\alpha) \leq T$, влизат в 
  съответните входни позиции на ОМ.
  
  \item (Б02) Съставя се абстрактен преход за ОМ (първоначално празен).
  
  \item (Б03) Проверява се дали текущото време е по-малко или равно на $T + t^{*}$.
  
  \item (Б04) Ако отговорът от (Б03) е „не“, се преминава към (Б12). В противен случай 
  се продължава с (Б05).
  
  \item (Б05) Определят се всички преходи, за които $t_1$ е по-голямо или равно на текущото време.
  
  \item (Б06) Типовете на намерените преходи от (Б05) се проверяват по следния начин:
  \begin{itemize}[itemsep=0.5em]
    \item (Б06а) Идентификаторите на входните позиции (участващи в типа на прехода) се заменят с 0, ако позицията е празна, и с 1, ако съдържа ядра.
    \item (Б06б) Изчислява се логическата стойност на получения булев израз.
  \end{itemize}

  \item (Б07) Всички преходи, които получават стойност 1 на (Б06), се добавят към абстрактния преход.
  
  \item (Б08) Прилага се алгоритъм {\bf A} към абстрактния преход.
  
  \item (Б09) От абстрактния преход се премахват всички преходи, които се деактивират в текущия момент.
  
  \item (Б10) Текущото време се увеличава с $t^{0}$.
  
  \item (Б11) Преминава се на стъпка (Б03).
  
  \item (Б12) ОМ преустановява функционирането си.
\end{itemize}

Използвайки двата алгоритъма, в \cite{Z1} е доказана следната теорема:


\begin{Theorem}\label{thm:th}
	В ОМ, която притежава всички свои компоненти, не възникват конфликтни ситуации.
\end{Theorem}



\subsection{Текущ софтуер за симулация на обобщени мрежи}



\subsubsection{Анализ на съществуващите софтуерни решения}

\begin{itemize}
  \item  \textbf{\textit{GN Lite}} \\
Това е пакет за моделиране и симулация на ОМ, изграден от независими модули и организиран в клиент--сървърна архитектура. Името \textit{GN Lite} се използва от 2003 г., когато се появява първият му компонент -- \textit{GNTicker}, високоефективен интерпретатор на ОМ модели \cite{Trifonov2005}. Платформата е проектирана с ясни стандартизирани точки на свързване: унифициран формат за описание на моделите и комуникационни протоколи към основния симулатор. Така всеки модул се разглежда като „черна кутия“, която може да бъде заменена с друг модул, стига да спазва същите спецификации.
Тази стандартна договореност позволява паралелна разработка на компонентите, използване на различни технологии от отделни екипи и безпроблемна интеграция на наследени инструменти (например средата \textit{Gennete}) като добавки. В резултат \textit{GN Lite} обединява по-широк набор от средства за работа с ОМ спрямо предшествениците си и е достигнал по-зрял етап на развитие. Планирани са още по-бързи алгоритми и добавяне на нови идеи от теорията на ОМ в следващи версии.

Клиент--сървърният модел осигурява гъвкавост при разполагане: симулационният сървър и потребителският интерфейс могат да се изпълняват на различни машини и операционни системи, комуникирайки през мрежов протокол. По-долу са описани основните подсистеми на \textit{GN Lite}, тяхното предназначение, използваните технологии, начините на взаимодействие и актуалното им състояние на разработка.


Ключови компоненти и функционалности на пакета:
\begin{itemize}[itemsep=0.5em]
	\item \textbf{\textit{GNTicker}}\\
	Това е централният симулатор на ОМ: зарежда модели във формат \textit{XGN}, изгражда позиции и преходи и управлява изпълнението на мрежата. Архитектурата отделя синтактичния разбор, планирането и проследяването, което улеснява разширяването и поддръжката \cite{Trifonov2005}. Вътрешно се използват индексирани матрици на инцидентност, както и детерминирана симулация с фиксиран начален параметър. Поддържани режими: стъпково, до условие и до краен хоризонт. Налични са проследяване и диагностика (например непостижими преходи, взаимни блокировки). Оптимизациите включват кеширане на резултатите от проверките за разрешимост и компактни структури за маркировки \cite{Dimitrov2011}. Предлагат се варианти за \textit{Win32} и \textit{GNU/Linux}.

	
	
	\item \textbf{\textit{XGN}}\\
	Това е междинен формат за обмен на модели, изграден върху \textit{XML} и проверяван спрямо \textit{XSD}. Спецификациите на \textit{W3C} за \textit{XML}~1.0 и \textit{XML Schema} гарантират еднозначна структура и автоматична проверка при четене и запис \cite{W3CXML,W3CXSD}. Моделната логика е отделена от данните за изглед (координати, оформление), поради което промени във визуализацията не променят семантиката. Пространствата от имена разграничават елементите на модела от визуалните метаданни, а стабилните идентификатори и постоянният ред на запис улесняват проследяването на версии и сравнения. Еволюцията на формата се подпомага чрез версии на схемата и преобразувания с \textit{XSLT} при запазена съвместимост \cite{Trifonov2008}.
	За разлика от ранни затворени формати, ограничавали междинната работа между инструменти, \textit{XGN} въвежда отворен, текстов и разширяем подход. Това носи: (i) независимост от платформа и език за програмиране; (ii) широка наличност на безплатни средства за четене, запис, преобразуване и обработка на \textit{XML}; (iii) естествено разширяване -- първоначално се описва структурата на мрежата, а впоследствие могат да се добавят координати и друга визуална информация; (iv) човешка четимост -- файловете могат да бъдат редактирани и с обикновен текстов редактор.
	
	По съдържание \textit{XGN} е организиран в четири части: преходи, позиции, ядра и функции. Разделът за функции е отворен -- не налага конкретен скриптов език и позволява избор според нуждите на модела.
	
	
	
	\item \textbf{Езиков слой: \textit{GNTCFL} и \textit{GNJS}}\\
    \textit{GNTCFL} е език за описване на условия (предикати) и характеристични функции в моделите \cite{Trifonov2008}. \textit{GNJS} е алтернативна реализация със синтаксис на \textit{JavaScript}, което улеснява писането на логика на познат език \cite{Trifonov2008}. \textit{JavaScript} е един от най-популярните програмни езици, което е допълнителна мотивация за такъв синтаксис \cite{LangPop}.	

	Езиковият слой предоставя:
	\begin{itemize}[itemsep=0.25em]
		\item готови операции за аритметика, множества и време;
		\item изолирана среда за изпълнение с ясно определени видове данни на вход/изход към ядрото;
		\item предварителна проверка на изразите и кеширане на резултати при изпълнение за по-висока скорост;
		\item отчетливи съобщения при грешки с посочване на мястото в изходния текст \cite{Trifonov2008}.
	\end{itemize}	
		
	
	\item \textbf{\textit{TickerServer} и \textit{GNTP}}\\
	\textit{TickerServer} предоставя \textit{GNTicker} като мрежова услуга и позволява едновременно изпълнение на множество симулации. Общуването с клиентските програми става чрез протокола \textit{GNTP} \cite{GNTickerProtocol}, чиито правила са близки до тези на \textit{HTTP} (като концепция за команди/отговори и кодове за състояние) \cite{Trifonov2008}. При работа под \textit{Unix}-подобни системи е възможно използване на съвместимостни слоеве за \textit{.NET} приложения, например \textit{Mono} \cite{Mono}.
	
    \textit{GNTP} предоставя: (i) постъпково изпълнение на модел; (ii) изпълнение до настъпване на зададено условие; (iii) подаване на ядра от клиента към изпълнителя; (iv) управление и съобщаване на грешки \cite{Trifonov2008,GNTickerProtocol}.
	
	Конфигурацията включва приемане на шифровани връзки (\textit{TLS}) на входната точка, проста проверка на достъп (потребител/парола или ядра за достъп) и при нужда разрешаване на междудоменен достъп за уеб клиенти \cite{Dimitrov2011,GNTickerProtocol}.


		
	\item \textbf{\textit{GN IDE}}\\
	Това е платформено независима среда на \textit{Java} \cite{JLS25} за моделиране и симулация, която общува с \textit{GNTicker} през \textit{GNTP} \cite{GNIDE2009}. Тя обединява редактиране на мрежи, проверка спрямо \textit{XSD}, визуално разполагане, наблюдение на изпълнението и управление на експерименти.
	В \textit{GN IDE} са реализирани оператори за редуциране (за оптимизация и свиване на модели при запазване на коректността) \cite{ReducingOps2015_2016}. Реализирани са и йерархични оператори за разгъване/работа с йерархии в моделите, които подпомагат повторната употреба на подмрежи и движение между нива на абстракция \cite{HierOps2022}.
	
	\item \textbf{\textit{GNVis}}\\	
	Този инструмент автоматично изчислява координати за елементите на \textit{XGN}-модели, за да осигури четливо начертание при липса на визуални данни в файла. Визуалната информация (координати, цветове) не е част от дефиницията на ОМ и не влияе на симулацията; \textit{GNTicker} я игнорира, затова тя трябва да се генерира отделно при нужда от изобразяване \cite{Trifonov2008}.
	
	\textit{GNVis} е конзолно приложение на \textit{C\#} за платформата \textit{.NET} и използва евристични методи за подреждане, целящи липса на застъпвания и минимален брой пресичания на ребра. Описани са както подходи на базата на „модел на сили“ \cite{Kamada1989}, така и генетични алгоритми за намиране на удобна за възприемане схема \cite{Trifonov2008}. Получените координати се записват обратно в \textit{XGN}, което ускорява последващите ръчни корекции.
	
	\textbf{Кога се използва:}
	\begin{itemize}[itemsep=0.25em]
		\item при отваряне на \textit{XGN} без визуална информация (първоначално разполагане);
		\item при нужда от автоматично повторно подреждане след многократни редакции (намаляване на застъпвания и пресичания).
	\end{itemize}
	
	Интегрирането към графична среда за моделиране е праволинейно (извикване на инструмента и включване на генерираните координати). Ограничение е зависимостта от \textit{Windows/.NET}; за други платформи е необходим слой за съвместимост като \textit{Mono} \cite{Mono}.
	
	\item \textbf{\textit{XGN2SVG}}\\
	Този компонент преобразува \textit{XGN} в \textit{SVG}, като използва \textit{XSLT} върху \textit{XML} представянето. Това осигурява платформена независимост и позволява визуализация директно в уеб браузър, без нужда от интегрирана среда. \textit{SVG} е векторен и мащабируем формат, удобен за редактиране и публикация. Генерираният \textit{SVG} може да включва стилове, пояснения и идентификатори, съвпадащи с \textit{XGN}-идентификаторите, както и сценарии за динамична анотация или проиграване на записите от проследяването. При необходимост от алтернативни разполагания могат да се използват инструменти за графи като \textit{Graphviz} \cite{Trifonov2008,SVG_F,Ellson2004}.


		
	\item \textbf{\textit{GNProfy}}\\
	Компонентът \textit{GNProfy} реализира обратната посока на трансформация: превръща \textit{UML} диаграми на последователности (\textit{XMI}) в изпълними модели на ОМ \cite{Trifonov2008,Koycheva2002}. Линиите на живот се тълкуват като позиции, събитията -- като преходи, а съобщенията -- като ребрени зависимости с времеви ограничения. Извършват се проверки за съгласуваност и се попълват липсващи параметри, след което се генерира \textit{XGN}, готов за симулация; разклоненията и повторенията се свеждат до подмрежи с еквивалентна семантика.
	
	\item \textbf{\textit{GNDraw} (настолен инструмент)}\\
	Това е лек инструмент на \textit{Java} за създаване и базова симулация на ОМ. Подходящ е за малки модели и бързи първоначални варианти; работи локално, без необходимост от сървърна част. Предлага бърз износ към \textit{XGN} и последващо отваряне в \textit{GN IDE} за разширени анализи и сценарии \cite{GNDraw}.
\end{itemize}

\item \textbf{\textit{GoGN}}\\
	Това е библиотека на \textit{C\#} за разработване на модели, при която се наследяват абстрактни класове и се компилират конкретни реализации. Подходът е насочен към програмисти и улеснява вграждането на ОМ в по-големи системи \cite{Gochev2012}. Изборът на \textit{C\#} и \textit{.NET} повишава производителността на разработката, но увеличава зависимостите от платформата; за \textit{Unix}-подобни системи са необходими междинни слоеве за съвместимост като \textit{Mono} \cite{Mono}. Предоставени са сигурни по типове програмни интерфейси за дефиниране на преходи и позиции, налични са описателни атрибути за метаданни и двойници за изпитания на модулно ниво; моделите могат да се изнасят към \textit{XGN} за съвместимост с останалите инструменти \cite{Gochev2012}.
\end{itemize}


\subsubsection{Основни недостатъци на текущите симулатори}

\hspace*{2.5em}Повечето текущи симулатори са реализирани с гъвкави архитектури и решения \cite{Dimitrov2010SoftwareGN}. Въпреки това имат съществени ограничения, които засягат едновременно архитектурата, потребителското изживяване и поддръжката.

\begin{itemize}
    \item \textbf{Фрагментираност и непълна интеграция.} \\
    Ранните версии на основния интерпретатор \textit{GNTicker} са конзолни и не предлагат стъпково/визуално проследяване, което затруднява интерактивния анализ на динамиката. По-късно необходимостта от интерактивност се компенсира чрез клиент--сървър архитектура с \textit{TickerServer} и протокола \textit{GNTP}, но това добавя допълнителни зависимости и оперативна сложност. Липсата на ранно интегрирана графична среда води до използването на отделни помощни инструменти, които функционират извън основния цикъл на моделиране и симулация.

    \item \textbf{Проблеми с преносимост и зависимости.} \\
    Част от инструментите са реализирани върху \textit{.NET/C\#} и първоначално са насочени към \textit{Windows}. Под \textit{Unix}-подобни ОС изпълнението често разчита на съвместимостни слоеве като \textit{Mono}, което понижава предвидимостта и увеличава усилията за инсталация и поддръжка. Технологичната хетерогенност затруднява изграждането на единен процес на разработка, тестване и внедряване (\textit{release}).

    \item \textbf{Езиков слой и съвместимост.} \\
    \textit{GNTicker} използва \textit{GNTCFL} за предикати и характеристични функции, но нишовият синтаксис увеличава кривата на обучение. Експерименталният вариант \textit{GNJS} цели да снижи този праг, но ранните му прототипи не осигуряват пълна съвместимост с \textit{GNTP}, ограничавайки дистанционния контрол и интеграцията със сървърната страна. В по-нов \textit{GNTicker} тази празнина е адресирана чрез съвместна поддръжка на \textit{JavaScript} и \textit{GNTP}, но историческият разрив е довел до фрагментирани потоци на работа.

    \item \textbf{Мащабируемост и производителност.} \\
    Алгоритмите за пренасяне на ядра в ранните реализации са коректни, но не винаги оптимални за много големи или силно свързани мрежи. Това се проявява в по-дълго време за изпълнение и по-висока ресурсоемкост при сложни сценарии. Докато оптимизациите не бъдат системно внедрени, използването за ресурсоемки модели изисква допълнително оразмеряване на хардуер и внимателно настройване.

    \item \textbf{Ограничения на инструментариума и екосистемата.}\\ 
    Инструменти като \textit{XGN2SVG} улесняват споделянето на резултати през уеб, но генерираният изход е статичен и не покрива нуждите от интерактивен дебъг и постъпково разглеждане. \textit{GNVis} предлага авторазполагане, но като отделен \textit{.NET} модул добавя платформени зависимости и не е плътно интегриран в цикъла редакция--симулация. \textit{GNProfy} изгражда мост към \textit{UML} екосистемата, но обхватът му е фокусиран основно върху диаграми на последователности и по-стари \textit{XMI} версии, което лимитира универсалността. Експерименталната рамка \textit{GoGN} предоставя мощен \textit{API} на \textit{C\#/.NET} и поддържа експорт към \textit{XGN} за съвместимост, но изисква програмиране и компилация, вместо декларативно моделиране, което повишава прага за навлизане.

\end{itemize}

    Текущото състояние на \textit{GN~Lite} демонстрира архитектурна фрагментираност, платформени зависимости и исторически разминавания в езиковата и протоколна съвместимост. Това се отразява в по-висок праг за усвояване, ограничена интерактивност, затруднена преносимост и непълна интеграция между редакция, визуализация и симулация. Нужна е еволюция към по-единна, платформено независима и потребителски ориентирана среда, която да запази силните страни на \textit{GN~Lite}, но да преодолее установените ограничения чрез системни оптимизации и интеграция около \textit{XGN}/\textit{GNTP}.
	
	Тези наблюдения мотивират разработването на \textit{OnlineGN} като уеб базирана, интегрирана и платформено независима среда, която обединява моделирането, визуализацията, редакцията и симулацията в единен работен процес.

\subsection{Цел и задачи на дисертационния труд}

\subsubsection{Формулиране на целта}
\hspace*{2.5em}Целта на настоящия дисертационен труд е да се разработи уеб базиран онлайн симулатор за ОМ -- \textit{OnlineGN}, който предоставя интерактивна среда за създаване, визуализация, редактиране, споделяне и симулация на модели, и да се валидира приложимостта му чрез разработване, симулация и верификация на реални ОМ модели.

\subsubsection{Задачи за постигане на целта}

\begin{enumerate}[
	label=\textbf{Задача \arabic*.},
	ref=\textbf{Задача \arabic*},   
	leftmargin=2.5em,
	labelsep=0.8em,
	align=left
	]
\item Разработване и имплементация на алгоритъм за автоматично изчертаване на ОМ по тяхно описание в структуриран текстов формат.
 \label{task:auto-draw}
 
 \item Разработване и имплементация на алгоритъм за преобразуване на \textit{SVG} изображение на ОМ в \textit{TeX} формат, с цел лесно включване на визуализации в научни публикации.
 \label{task:svg-to-tex}
 
 \item Формулиране на изисквания към уеб базирано решение.
 \label{task:requirements}
 
 \item Проектиране на архитектурата на системата \textit{OnlineGN} и дефиниране на основните модули (съхранение на мрежи, визуализация, редакция, симулация, импорт/експорт, \textit{API}).
 \label{task:architecture}
 
 \item Реализация на уеб базиран графичен интерфейс за визуализация, редакция и персонализация на генерираното изображение (позиции, преходи, връзки и визуални параметри).
 \label{task:web-ui}
 
 \item Реализация на удобен интерфейс за редакция на предикати в индексираната матрица и характеристични функции на позициите.
 \label{task:predicates-ui}
 
 \item Реализация на механизъм за онлайн съхранение и споделяне на мрежи чрез уникална хипервръзка (линк), базиран на централизирано съхранение в база данни.
 \label{task:share-storage}
 
  \item Създаване и формализиране на нови ОМ модели на реални процеси, които до момента не са описвани чрез ОМ.
 \label{task:new-models}
 
  \item Симулация и верификация на разработените модели в \textit{OnlineGN} чрез сценарии и експерименти.
  \label{task:simulation}
\end{enumerate}


Реализирането на поставените задачи ще доведе до създаването на функционален уеб базиран симулатор, придружен от нови алгоритми за автоматизирано изчертаване и преобразуване към \textit{TeX}. Практическата приложимост на разработката ще бъде потвърдена чрез валидиращи примери и формализиране на нови ОМ модели на реални процеси.

\vspace{0.5cm}

В обобщение, дисертационният труд адресира необходимостта от автоматизация и ефективно сътрудничество при работа с апарата на ОМ. Проектът \textit{OnlineGN} предлага цялостно, гъвкаво и лесно за употреба решение, съобразено със съвременните софтуерни стандарти и нуждите на научноизследователската общност.

%\end{enumerate}